\chapter{One-dimensional continuum mechanics}
\label{ch:continuum}

\checkedagainst{1.9.0 (\code{8011ac3})}{The derivations are checked
symbolically in \code{doc/manual/checks/ch2_continuum.py} (sympy, for any
smooth flow); the units by arithmetic in the same script; what the code does
at $t = 0$ against \code{maincode.f95}, section~5.}

This chapter sets up the equations that every later chapter solves. It
starts from the balance laws of continuum mechanics, narrows them to the
one regime \bw\ computes (planar flow in uniaxial strain, with no material
strength), and rewrites them in the Lagrangian form the code is built on:
in a coordinate that moves with the material and counts mass instead of
distance. It ends with what the equations say about entropy, small
disturbances, the edges of a problem, and how \bw\ lays its mesh over a
stack of layers at $t = 0$.

The development follows the ME~697 course (its Units~3 and~4); the
notation is this manual's (\cref{app:notation}).

\section{Uniaxial strain}
\label{sec:uniaxial}

A plate-impact experiment is built so that, for a while, the material at
its center moves only along one axis. The flyer and the target are wide
compared with their thickness, so waves from the edges take time to reach
the center, and until they do the motion there is planar: everything
depends on one coordinate, $x$, and on time, $t$. This is \emph{uniaxial
strain}: the material is compressed or stretched along $x$ and not at all
across it (\cref{fig:prism}).

\begin{figure}[ht]
\centering
\includegraphics[width=0.5\textwidth]{figures/prism_uniaxial.pdf}
\caption{Uniaxial strain. A prism of material (dashed) is compressed along
$x$ alone (solid): its length changes, its cross-section does not.}
\label{fig:prism}
\end{figure}

In uniaxial strain:

\begin{itemize}
\item the velocity has one nonzero component, $\langle u, v, w\rangle =
  \langle u, 0, 0\rangle$, and every field depends on $(x, t)$ alone:
  $\partial/\partial y = \partial/\partial z = 0$;
\item there is no shear and no rotation, so the stress tensor is diagonal,
  \begin{equation}
  \bm{\sigma} = \begin{pmatrix}
    \sigma_{xx} & 0 & 0 \\ 0 & \sigma_{yy} & 0 \\ 0 & 0 & \sigma_{zz}
  \end{pmatrix},
  \label{eq:cm-stress}
  \end{equation}
  and only $\sigma_{xx}$, the stress on a face normal to $x$, enters the
  equations of motion. It is positive in tension, as is usual in continuum
  mechanics; pressure $P$ is positive in compression;
\item diffusion (heat conduction, and the diffusion of species) is left
  out, because it is too slow to matter. Heat diffuses across a length $L$
  in a time of order $L^2/\alpha$, where $\alpha$ is the thermal
  diffusivity. For example, it takes about 10~ms to cross one millimeter of
  aluminum and about 10~s to cross one millimeter of a plastic, whereas a
  plate-impact run ends within a few microseconds, three to six orders of
  magnitude sooner. Even at a shock front, where the temperature changes
  most sharply, conduction acts only over a width of order $\alpha/U_s$;
  this is about 15~nm in aluminum and far less in a plastic, which is far
  below the size of any cell in a mesh.
\end{itemize}

\section{Stress, and why \bw\ is hydrodynamic}
\label{sec:hydro-only}

The balance laws below hold for any material; what distinguishes one
material from another is how $\sigma_{xx}$ depends on its state. The
simplest choices are:

\begin{center}\small
\begin{tabularx}{\textwidth}{@{}l l Y@{}}
\toprule
Material & $\sigma_{xx}$ & \\
\midrule
ideal gas or fluid & $-P$ & no resistance to shear \\
Newtonian fluid & $-P + \tfrac{4}{3}\mu\,\partial u/\partial x$ & viscosity
  $\mu$; the bulk viscosity taken as zero (Stokes' hypothesis) \\
yielding solid & $-P - \tfrac{2}{3}Y$ & von Mises yield, strength $Y$ \\
\bw & $-(P + q)$ & no strength; $q$ is method AV's artificial viscosity
  (\cref{ch:method-av}), zero in method RS \\
\bottomrule
\end{tabularx}
\end{center}

\bw\ is \emph{hydrodynamic}: its materials carry pressure but no
strength, $Y = 0$. This is a choice, not a missing piece. Above a few GPa,
the pressures of shock and detonation physics, a solid's strength is a
small part of its stress, and a hydrocode with a good equation of state
captures the waves; leaving strength out keeps the code compact and every
equation in it checkable.
% SCOPE: strength (out of scope for good, DK, Nov 19, 2026).
\bw\ will not gain a strength model.\footnote{A
decision of the project (\code{TODO.md}, Priorities, Nov 19, 2026). The
deck key \key{spall_strength} is not a strength model: it limits how much
tension the pressure may hold before a cell breaks (\cref{ch:materials}).}
So throughout this manual
\begin{equation}
\sigma_{xx} = -(P + q).
\label{eq:cm-sigma}
\end{equation}

\section{The balance laws}
\label{sec:balance}

\paragraph{A scalar carried by the flow.} Let $\phi$ be an amount of
something per unit volume (mass, momentum, energy), and $\tilde\phi =
\phi/\rho$ the same amount per unit mass. In one dimension, the
\emph{conservative} form of the balance law is given by
\begin{equation}
\frac{\partial \phi}{\partial t} + \frac{\partial}{\partial x}(u\phi)
  = S_\phi,
\label{eq:cm-transport}
\end{equation}
The amount of $\phi$ in a fixed slab changes in two ways: the flow carries it in or out across the slab's faces, at the
rate $u\phi$, and it is produced or consumed within the slab, at the rate
$S_\phi$ per unit volume (the \emph{source}). Writing $\phi =
\rho\tilde\phi$, and using conservation of mass (the same balance with
$\phi = \rho$ and no source, below) to remove the derivatives of $\rho$,
gives the \emph{intensive} form,
\begin{equation}
\rho\,\frac{D\tilde\phi}{Dt} = S_\phi,
\qquad
\frac{D}{Dt} = \frac{\partial}{\partial t} + u\,\frac{\partial}{\partial x},
\label{eq:cm-intensive}
\end{equation}
in which $D/Dt$ is the \emph{material derivative}: the rate of change seen
by an observer riding with the material.

\paragraph{A slab that moves.} The same balance can be written for a slab
whose faces move. Integrate \cref{eq:cm-transport} from a left face at
$x_a(t)$ to a right face at $x_b(t)$; Leibniz's rule for an integral with
moving limits gives
\begin{equation}
\frac{\dd}{\dd t}\int_{x_a(t)}^{x_b(t)} \phi\,\dd\hat x
  = -\Bigl[\,\phi\,\Bigl(u - \frac{\dd x}{\dd t}\Bigr)\Bigr]_{x_a}^{x_b}
    + \int_{x_a(t)}^{x_b(t)} S_\phi\,\dd\hat x .
\label{eq:cm-rtt}
\end{equation}
This is the Reynolds transport theorem in one dimension. The bracket is the
flux \emph{relative to the faces}. If the faces move with the material,
$\dd x/\dd t = u$, nothing crosses them: the slab keeps the same matter for
all time. That is a Lagrangian cell, and it is why \bw\ moves its mesh with
the material: the advection terms vanish, and only the sources remain.
(\Cref{ch:shocks} uses the same theorem the other way, with faces fixed
to a moving shock.)

\paragraph{The sources.} A slab of material between $x$ and $x + \Delta x$
(\cref{fig:fbd}) feels the stress $\sigma_{xx}$ on its left face and
$\sigma_{xx} + \Delta\sigma_{xx}$ on its right. Per unit volume, the net
force is the momentum source, and the work those stresses do at the face
velocities $u$ and $u + \Delta u$ is the energy source:
\begin{equation}
S_{\mathrm{mom}} = \frac{\partial \sigma_{xx}}{\partial x},
\qquad
S_{\mathrm{tot}} = \frac{\partial}{\partial x}(\sigma_{xx} u)
  = \underbrace{u\,\frac{\partial \sigma_{xx}}{\partial x}}_{\text{kinetic}}
  + \underbrace{\sigma_{xx}\frac{\partial u}{\partial x}}_{\text{internal}}.
\label{eq:cm-sources}
\end{equation}
The total work splits in two. The first part accelerates the slab: it is
the momentum source times $u$, the source of kinetic energy. The second,
the \emph{stress power}, deforms it: the source of internal energy.

\begin{figure}[ht]
\centering
\includegraphics[width=0.55\textwidth]{figures/fbd_uniaxial.pdf}
\caption{A small element of material in uniaxial strain (solid, liquid or
gas). The stress on its two faces normal to $x$ differs by
$\Delta\sigma_{xx}$, and the velocity by $\Delta u$. $\sigma_{yy}$ acts on
the other faces but does no work, since nothing moves along $y$ or $z$.}
\label{fig:fbd}
\end{figure}

\paragraph{Mass, momentum and energy.} With $\phi = \rho$, $\rho u$ and
$\rho e$ in turn, and the sources above, the balance laws in intensive
form are
\begin{align}
\frac{D\rho}{Dt} &= -\rho\,\frac{\partial u}{\partial x},
\label{eq:cm-mass}\\
\rho\,\frac{Du}{Dt} &= \frac{\partial \sigma_{xx}}{\partial x},
\label{eq:cm-mom}\\
\rho\,\frac{De}{Dt} &= \sigma_{xx}\,\frac{\partial u}{\partial x},
\label{eq:cm-energy}
\end{align}
with $e$ the specific internal energy. Their sum (\cref{eq:cm-mom} times
$u$, plus \cref{eq:cm-energy}) is the balance of total energy, $E = e +
\tfrac{1}{2}u^2$:
\begin{equation}
\rho\,\frac{DE}{Dt} = \frac{\partial}{\partial x}(\sigma_{xx} u).
\label{eq:cm-total}
\end{equation}
For a fluid with $\sigma_{xx} = -P$ these are the compressible Euler
equations in one dimension. A reaction adds nothing to them in \bw: a
reacting cell's chemistry is inside its equation of state, through the
closure of \cref{ch:burning-cell}. There $e$ is the whole internal energy,
chemical energy included, and a reaction turns part of it into heat
without changing it.\footnote{The course's sheets carry the heat of reaction as an
explicit source, $\sum_k Q_k\rho\dot\Omega_k$, with $e$ thermal only. The
two bookkeepings describe the same physics; \bw\ uses the second, so its
total energy is conserved as written in \cref{eq:cm-total}.}

\section{The mass coordinate}
\label{sec:mass-coordinate}

A Lagrangian code labels each piece of material once and then follows it
for all time. The
natural label in one dimension is the mass to its left: the \emph{areal
mass coordinate}
\begin{equation}
h(x, t) = \int_{x_L(t)}^{x} \rho(\hat x, t)\,\dd\hat x,
\label{eq:cm-h}
\end{equation}
mass per unit area, measured from the left edge of the material,
$x_L(t)$. (It is written $h$ as in the course; it is not an enthalpy, and
no enthalpy appears in this manual.) Its derivatives are
\begin{equation}
\frac{\partial h}{\partial x} = \rho,
\qquad
\frac{\partial h}{\partial t}\Big|_x = -\rho u,
\label{eq:cm-hderiv}
\end{equation}
where the second expression shows that the mass to the left of a fixed
point changes only by
the flux $\rho u$ across it (the two are compatible exactly when mass is
conserved, \cref{eq:cm-mass}). A material point keeps its $h$ for all
time, so in the coordinates $(h, t)$:
\begin{equation}
\frac{\partial}{\partial t}\Big|_h = \frac{D}{Dt},
\qquad
\frac{\partial}{\partial h} = \frac{1}{\rho}\frac{\partial}{\partial x}
  = V\frac{\partial}{\partial x},
\label{eq:cm-hrules}
\end{equation}
with $V = 1/\rho$ the specific volume. The time derivative at fixed $h$
is the material derivative, and the space derivative is per unit of mass
rather than of length.

\paragraph{The equations in $h$.} Substituting \cref{eq:cm-hrules} into
\cref{eq:cm-mass,eq:cm-mom,eq:cm-energy,eq:cm-total}, and using
$D\rho/Dt = -\rho^2\,DV/Dt$ in the first, gives the form \bw\ is built on:
\begin{align}
\frac{\partial V}{\partial t} &= \frac{\partial u}{\partial h},
\label{eq:lag-mass}\\
\frac{\partial u}{\partial t} &= \frac{\partial \sigma_{xx}}{\partial h},
\label{eq:lag-mom}\\
\frac{\partial e}{\partial t} &= \sigma_{xx}\,\frac{\partial V}{\partial t}
  = -(P + q)\,\frac{\partial V}{\partial t},
\label{eq:lag-energy}\\
\frac{\partial}{\partial t}\Bigl(e + \tfrac{1}{2}u^2\Bigr)
  &= \frac{\partial}{\partial h}(\sigma_{xx} u),
\label{eq:lag-total}
\end{align}
where every time derivative corresponds to a fixed value of $h$. The
density never needs an equation of
its own: it follows from where the material is,
$\partial x/\partial h = V$. And the energy equation has become the first
law for a piece of material, $\dd e = -P\,\dd V - q\,\dd V$: its energy
changes by the work of the pressure and of the viscosity as its volume
changes, and by nothing else.

\paragraph{A cell.} Integrate \cref{eq:lag-mass,eq:lag-mom,eq:lag-total}
over a slab between two material points $h_a$ and $h_b$. The slab's mass
per unit area, $M = h_b - h_a$, is fixed for all time; its length, its
momentum and its total energy change only through its faces:
\begin{equation}
\frac{\dd}{\dd t}(x_b - x_a) = u_b - u_a,
\quad
M\,\frac{\dd \bar u}{\dd t} = \sigma_b - \sigma_a,
\quad
M\,\frac{\dd \bar E}{\dd t} = \sigma_b u_b - \sigma_a u_a,
\label{eq:cm-cell}
\end{equation}
with $\bar u$ and $\bar E$ the slab's mean velocity and total energy per
unit mass, and $\sigma_a$, $u_a$ ($\sigma_b$, $u_b$) the stress and
velocity at its left (right) face. These are exact. A face shared by two
cells appears in both with opposite signs, so in a sum over every cell all
the inner faces cancel: the whole stack's momentum and energy change only
through its two ends. \Cref{ch:method-av,ch:method-rs} are two ways of
deciding what $u$ and $\sigma$ are at a face, between time levels; both
solution methods keep this cancellation, which is why both conserve momentum and total
energy to round-off.

\section{Entropy and the second law}
\label{sec:entropy}

For a material whose state is set by two variables, the Gibbs relation
$T\,\dd s = \dd e + P\,\dd V$ defines the specific entropy $s$. Following a
piece of material, and using \cref{eq:lag-energy}:
\begin{equation}
T\,\frac{\partial s}{\partial t}
  = \frac{\partial e}{\partial t} + P\,\frac{\partial V}{\partial t}
  = (\sigma_{xx} + P)\,\frac{\partial V}{\partial t}
  = -\,q\,\frac{\partial V}{\partial t},
\label{eq:cm-entropy}
\end{equation}
or, in the Eulerian variables of the course, $\rho\,Ds/Dt = \bigl((\sigma_{xx}
+ P)/T\bigr)\,\partial u/\partial x$. Three consequences:

\begin{itemize}
\item \textbf{A fluid with $\sigma_{xx} = -P$ is isentropic wherever it is
  smooth.} Only the part of the stress that is not pressure can generate
  entropy. A non-linear shock does produce a rise in entropy, so the smooth
  equations alone cannot describe entropy production; \cref{ch:shocks}
  shows how a shock of finite thickness does.
\item \textbf{The second law constrains $q$.} Entropy must not fall,
  $\partial s/\partial t \ge 0$, so $q\,\partial V/\partial t \le 0$:
  wherever $q \ne 0$ the material must be compressed ($\partial
  V/\partial t < 0$) and $q > 0$. Method AV's viscosity is built to obey
  this: it is switched on only in a cell that is being compressed
  (\cref{ch:method-av}).
\item \textbf{A Newtonian fluid obeys it automatically.} With $\sigma_{xx}
  + P = \tfrac{4}{3}\mu\,\partial u/\partial x$, the right-hand side is
  $\tfrac{4}{3}(\mu/T)(\partial u/\partial x)^2 \ge 0$. Artificial
  viscosity borrows this structure (\cref{ch:shocks}).
\end{itemize}

\section{Small disturbances}
\label{sec:acoustics}

Disturb material at rest, $(\rhoz, P = 0, u = 0)$, a little, so that the
changes are small and isentropic. The pressure then follows the volume
through the isentropic sound speed $c$,
\begin{equation}
\dd P = -\rhoz^2 c^2\,\dd V,
\qquad
c^2 = -V^2\Bigl(\frac{\partial P}{\partial V}\Bigr)_s ,
\label{eq:cm-c}
\end{equation}
and \cref{eq:lag-mass,eq:lag-mom} (with $q = 0$) combine to the wave
equation in $h$,
\begin{equation}
\frac{\partial^2 u}{\partial t^2} = (\rhoz c)^2\,\frac{\partial^2 u}{\partial h^2},
\qquad
u = f(h - \rhoz c\,t) + g(h + \rhoz c\,t).
\label{eq:cm-wave}
\end{equation}
A disturbance travels through the material's mass at $\rhoz c$, which in
distance is the sound speed $c$, in either direction. The product
\begin{equation}
Z = \rhoz c
\label{eq:cm-impedance}
\end{equation}
is the \emph{acoustic impedance}, and it is also the ratio of the two
fields in a single wave: a right-running wave carries $P = Z u$, a
left-running one $P = -Z u$. Both facts reappear in the numerics. The
speed $c$ sets how long a step may be (no wave may cross more than one
cell in one step, \cref{ch:method-av,ch:method-rs}), and the impedance is
the weight method RS gives each side of a face (\cref{ch:method-rs}).
Strong shocks travel faster than $c$ and carry more pressure per unit of
velocity than $Z$; \cref{ch:shocks} replaces both with their shock
counterparts.

\section{Interfaces and boundaries}
\label{sec:boundaries}

\paragraph{Contacts.} \bw\ computes a stack of layers in contact: a
flyer, a target, a window. At a contact between two materials, the
balance of momentum across a slab of vanishing mass makes the stress
continuous, and the material on the two sides never separates or
interpenetrates while the contact holds, so the velocity is continuous
too:
\begin{equation}
u^- = u^+, \qquad \sigma_{xx}^- = \sigma_{xx}^+ .
\label{eq:cm-contact}
\end{equation}
Density, internal energy and temperature jump across a contact. \bw\
places every contact on a face of its mesh (below), so no cell ever holds
two materials.

\paragraph{Free surfaces.} The two ends of the stack face a void, which
carries no stress:
\begin{equation}
\sigma_{xx} = 0 \quad\text{at both ends.}
\label{eq:cm-free}
\end{equation}
These are the only boundaries \bw\ has: no rigid wall, no piston, no
inflow. A layer at rest is at zero pressure, so its free surface is in
equilibrium at $t = 0$. That is why every material in \bw\ starts at
$P = 0$: the solids and liquids of shock physics, whose ambient pressure
is negligible next to the pressures of interest.
% SCOPE: gases. Update this paragraph (and the boundary list above) when
% bitwhiz gains a gas EOS or new boundaries (doc/manual/README.md, Scope).
A gas at ambient pressure, whose ambient state is the scale of its shocks,
is outside the present scope.

\paragraph{Breaking.} A contact can also open. When the material near it
is pulled harder than it can hold, it breaks (spalls), and the two sides
then behave as free surfaces until they collide again. How much tension a
material holds, and what a broken cell does, is part of its material model
(\cref{ch:materials}).

\section{The mesh at the start of a run}
\label{sec:mesh}

A deck describes a stack of layers, left to right, each with a material,
a thickness and a number of cells (or a resolution, in cells per mm), and
a starting velocity. At $t = 0$ \bw\ builds the mesh from them
(\code{maincode.f95}, section~5):

\begin{enumerate}
\item \textbf{Faces.} Each layer's faces are evenly spaced, face $j$ of a
  layer of $n$ cells at $x_{\mathrm{left}} + j\,w/n$, with $w$ the layer's
  thickness. Each position is computed from the layer's edge, not by
  adding up cell widths, so round-off cannot accumulate across a layer, and
  each layer's left edge, a contact, lands exactly on a face.
\item \textbf{Cells.} Cell $k$, between faces $k - \tfrac{1}{2}$ and
  $k + \tfrac{1}{2}$, takes its layer's material constants and the
  reference state
  \begin{equation}
  \rho_k = \rhoz, \quad e_k = 0, \quad P_k = 0, \quad q_k = 0,
  \quad c_k = C_0,
  \label{eq:cm-start}
  \end{equation}
  in which the equation of state must return $P = 0$ and $c = C_0$; the
  setup checks that it does. Its mass per unit area,
  \begin{equation}
  M_k = \rhoz\,\bigl(x_{k+1/2} - x_{k-1/2}\bigr),
  \label{eq:cm-mass-cell}
  \end{equation}
  is fixed from then on. Units: g/cm$^3$ times mm.
\item \textbf{Velocities.} Each cell starts at its layer's velocity, or,
  with a velocity profile (deck \key{xvel_profile}), at the profile's mean
  over the cell, chosen so that the layer's momentum is exactly its
  thickness times $\rhoz$ times its mean velocity. A face between cells
  of equal velocity takes that velocity. A face between cells of different
  velocities (an \emph{impact face}, where a flyer meets its target, or a
  face inside a profile) takes their momentum-weighted mean,
  \begin{equation}
  u_{k+1/2} = \frac{M_k u_k + M_{k+1} u_{k+1}}{M_k + M_{k+1}},
  \label{eq:cm-impact}
  \end{equation}
  so the face carries exactly the momentum of the two half cells it
  stands for, and the stack's momentum at $t = 0$ is the deck's.
  \trace{modules/dist_mod.f95}{dist_cdf, dist_sf}{test_dist}
\item \textbf{The ends.} Beyond each end of the stack sits a ghost cell
  of zero mass and zero pressure: the void of \cref{eq:cm-free}.
\end{enumerate}

\paragraph{Units.} \bw's units are millimeters, microseconds, g/cm$^3$,
km/s (= mm/\textmu s), GPa, kJ/g (= (km/s)$^2$) and kelvin. They form a
consistent set, so no equation in the code carries a conversion factor:
\begin{center}\small
\begin{tabular}{@{}l l@{}}
\toprule
Product & Units \\
\midrule
$\rho u^2$ & (g/cm$^3$)(km/s)$^2$ = GPa \\
$P V$ & GPa\,/\,(g/cm$^3$) = kJ/g \\
$\Delta t\,\Delta P / M$ & \textmu s\,GPa\,/\,(g/cm$^3$\,mm) = km/s \\
$\rho c$ & (g/cm$^3$)(km/s) = GPa per km/s \\
\bottomrule
\end{tabular}
\end{center}
The third line is the momentum update of a cell, \cref{eq:cm-cell}: a
pressure difference acting over a step, divided by a cell's mass, is a
change of velocity, with nothing to convert.
